% from ega1-10-fr.tex
\oldpage[I]{211}
  \item[\rm{(ii)}] Soient $\mathfrak{X}$, $\mathfrak{Y}$, $\mathfrak{S}$ trois
  préschémas formels localement noethériens,
  $f:\mathfrak{X}\to\mathfrak{S}$ une immersion fermée,
  $g:\mathfrak{Y}\to\mathfrak{S}$ un morphisme. Alors le morphisme
  $\mathfrak{X}\times_{\mathfrak{S}}\mathfrak{Y}\to\mathfrak{Y}$ est une
  immersion fermée.
  \item[\rm{(iii)}] Soient $\mathfrak{S}$ un préschéma formel localement
  noethérien, $\mathfrak{X}'$, $\mathfrak{Y}'$ deux
  $\mathfrak{S}$-préschémas formels localement noethériens tels que
  $\mathfrak{X}'\times_{\mathfrak{S}}\mathfrak{Y}'$ soit localement
  noethérien. Si $\mathfrak{X}$, $\mathfrak{Y}$ sont des
  $\mathfrak{S}$-préschémas formels localement noethériens,
  $f:\mathfrak{X}\to\mathfrak{X}'$, $g:\mathfrak{Y}\to\mathfrak{Y}'$ deux
  $\mathfrak{S}$-morphismes qui sont des immersions fermées, alors
  $f\times_{\mathfrak{S}}g$ est une immersion fermée.
\end{enumerate}
\end{proposition}

En vertu de (\hyperref[I.3.5.1-fr]{3.5.1}), il suffit encore de prouver (i)
et (ii).

Pour démontrer (i), on peut supposer que $\mathfrak{Y}$ (resp.
$\mathfrak{Z}$) est un sous-préschéma fermé de $\mathfrak{X}$ (resp.
$\mathfrak{Y}$) défini par un faisceau cohérent $\mathscr{J}$ (resp.
$\mathscr{K}$) d'idéaux de $\mathscr{O}_{\mathfrak{X}}$ (resp.
$\mathscr{O}_{\mathfrak{Y}}$)~; si $\psi$ est l'injection
$\mathfrak{Y}\to\mathfrak{X}$ des espaces sous-jacents,
$\psi_*(\mathscr{K})$ est un faisceau cohérent d'idéaux de
$\psi_*(\mathscr{O}_{\mathfrak{Y}})
=\mathscr{O}_{\mathfrak{X}}/\mathscr{J}$
(\hyperref[0.5.3.12-fr]{0, 5.3.12}), donc aussi un
$\mathscr{O}_{\mathfrak{X}}$-Module cohérent
(\hyperref[0.5.3.10-fr]{0, 5.3.10})~; le noyau $\mathscr{K}_1$ de
$\mathscr{O}_{\mathfrak{X}}\to
(\mathscr{O}_{\mathfrak{X}}/\mathscr{J})/\psi_*(\mathscr{K})$ est donc un
faisceau cohérent d'idéaux de $\mathscr{O}_{\mathfrak{X}}$
(\hyperref[0.5.3.4-fr]{0, 5.3.4}), et
$\mathscr{O}_{\mathfrak{X}}/\mathscr{K}_1$ est isomorphe à
$\psi_*(\mathscr{O}_{\mathfrak{Y}}/\mathscr{K})$, ce qui prouve que
$\mathfrak{Z}$ est isomorphe à un sous-préschéma fermé de $\mathfrak{X}$.

Pour démontrer (ii), il est immédiat qu'on peut se limiter au cas où
$\mathfrak{S}=\operatorname{Spf}(A)$,
$\mathfrak{X}=\operatorname{Spf}(B)$,
$\mathfrak{Y}=\operatorname{Spf}(C)$, $A$ étant un anneau
$\mathfrak{J}$-adique noethérien, $B=A/\mathfrak{a}$, où $\mathfrak{a}$ est
un idéal de $A$, $C$ une $A$-algèbre topologique adique et noethérienne. Tout
revient à prouver que l'homomorphisme
$C\to C\widehat{\otimes}_A(A/\mathfrak{a})$ est \emph{surjectif}~: or,
$A/\mathfrak{a}$ est un $A$-module de type fini, et sa topologie est la
topologie $\mathfrak{J}$-adique~; il résulte alors de
(\hyperref[0.7.7.8-fr]{0, 7.7.8}) que
$C\widehat{\otimes}_A(A/\mathfrak{a})$ s'identifie à
$C\otimes_A(A/\mathfrak{a})=C/\mathfrak{a}C$, d'où notre assertion.

\begin{corollary}[10.14.6]
\label{I.10.14.6-fr}
Sous les hypothèses de (\hyperref[I.10.14.5-fr]{10.14.5, (ii)}), soient
$p:\mathfrak{X}\times_{\mathfrak{S}}\mathfrak{Y}\to\mathfrak{X}$,
$q:\mathfrak{X}\times_{\mathfrak{S}}\mathfrak{Y}\to\mathfrak{Y}$ les
projections, de sorte que le diagramme
\[
  \xymatrix{
    \mathfrak{X} \ar[d]_{f}
    & \mathfrak{X}\times_{\mathfrak{S}}\mathfrak{Y}
      \ar[l]_{p} \ar[d]^{q}\\
    \mathfrak{S}
    & \mathfrak{Y} \ar[l]^{g}
  }
\]
est commutatif. Pour tout $\mathscr{O}_{\mathfrak{X}}$-Module cohérent
$\mathscr{F}$, on a alors un isomorphisme canonique de
$\mathscr{O}_{\mathfrak{Y}}$-Modules
\[
  \label{I.10.14.6.1-fr}
  u:g^*(f_*(\mathscr{F}))
  \xrightarrow{\sim}
  q_*(p^*(\mathscr{F})).
  \tag{10.14.6.1}
\]
\end{corollary}

Pour définir un homomorphisme
$g^*(f_*(\mathscr{F}))\to q_*(p^*(\mathscr{F}))$, on sait qu'il revient au
même de définir un homomorphisme
$f_*(\mathscr{F})\to g_*(q_*(p^*(\mathscr{F})))
=f_*(p_*(p^*(\mathscr{F})))$
(\hyperref[0.4.4.3-fr]{0, 4.4.3})~: nous prendrons $u=f_*(\rho)$, où $\rho$
est l'homomorphisme canonique
$\mathscr{F}\to p_*(p^*(\mathscr{F}))$
(\hyperref[0.4.4.3-fr]{0, 4.4.3}). Pour voir que $u$ est un isomorphisme, on
se ramène aussitôt au cas où $\mathfrak{S}$, $\mathfrak{X}$,
$\mathfrak{Y}$ sont des spectres formels d'anneaux adiques noethériens $A$,
$B$, $C$, avec les conditions vues ci-dessus dans
(\hyperref[I.10.14.5-fr]{10.14.5, (ii)})~; on a alors
$\mathscr{F}=M^\Delta$, où $M$ est un $(A/\mathfrak{a})$-module de type fini
(\hyperref[I.10.10.5-fr]{10.10.5}), et les deux membres de
(\hyperref[I.10.14.6.1-fr]{10.14.6.1}) s'identifient respectivement, en vertu
de (\hyperref[I.10.10.8-fr]{10.10.8}), à
$(C\otimes_A M)^\Delta$ et à
$((C/\mathfrak{a}C)\otimes_{A/\mathfrak{a}}M)^\Delta$, d'où le corollaire,
puisque
$(C/\mathfrak{a}C)\otimes_{A/\mathfrak{a}}M
=(C\otimes_A(A/\mathfrak{a}))\otimes_{A/\mathfrak{a}}M$ s'identifie
canoniquement à $C\otimes_A M$.

